\section{Framework Design}
% GOAL: Describe the evaluation framework — AIhub unified access, containerized isolation, workload design, resurrection protocol.

\subsection{Unified Evaluation Infrastructure (AIhub)}

\evobench{} is built on \textsc{AIhub}, a unified API gateway (\url{https://aihub.arcsysu.cn}) that provides standardized access to 21 large language models spanning both proprietary and open-weight families.
This design serves two purposes:
(1)~\emph{model diversity}: the framework can evaluate any model accessible through an OpenAI-compatible API, enabling systematic comparison across model families, scales, and training paradigms;
(2)~\emph{backend diversity}: the framework supports multiple agent scaffolds---OpenHands SDK (full agent framework with unlimited self-loop), Claude Code (file I/O + command execution), Codex CLI (autonomous coding with sandboxed execution), and Kimi CLI (conversational coding with plan management)---each representing a distinct interaction paradigm.

All models are accessed through a unified interface, eliminating API-specific integration overhead and ensuring that evaluation differences reflect model and scaffold capabilities rather than infrastructure variability.

\subsection{Containerized Isolation}

Each evaluation run operates in an isolated workspace containing a complete copy of the target repository with all dependencies (ANTLR, LLVM 14, pybind11) pre-installed.
For \yatcc{}-Hard, each run additionally launches a fresh container environment that is destroyed after completion, ensuring zero cross-run contamination.
The workspace is configured via CMake with a fixed toolchain; scoring scripts are deterministic, producing identical scores for identical code regardless of execution order or timing.

\subsection{Workload Design: Serial Dependency Chains}

\evobench{} provides two complementary workloads grounded in compiler construction, a domain chosen for its strict causal dependencies, well-defined intermediate artifacts, and objective automated scoring:

\subsubsection{\yatcc{}: Scaffolded Implementation}

The original \yatcc{} workload~\citep{yatcc2024} organizes six tasks along the compiler construction pipeline (Table~\ref{tab:tasks}).
Each task provides \emph{scaffolded starter code}---the agent fills in missing implementation details within a provided framework.
This workload measures an agent's ability to \emph{complete} a partially-implemented system.

\begin{table}[t]
\centering
\small
\begin{tabular}{clll}
\toprule
\textbf{Task} & \textbf{Name} & \textbf{Objective} & \textbf{Input Artifact} \\
\midrule
0 & Environment Setup & Build \& verify toolchain & --- \\
1 & Lexer & Tokenize C source code & Preprocessed source \\
2 & Parser & Construct AST/ASG & Token stream (from T1) \\
3 & IR Generation & Emit LLVM IR & AST/JSON (from T2) \\
4 & IR Optimization & Implement LLVM Pass & Raw IR (from T3) \\
5 & Code Generation & Produce RV64 assembly & Optimized IR (from T4) \\
\bottomrule
\end{tabular}
\caption{\evobench{} task chain. Each task consumes the output artifact of its predecessor, forming a strict serial dependency.}
\label{tab:tasks}
\end{table}

\subsubsection{\yatcc{}-Hard: Full From-Scratch Implementation}

\yatcc{}-Hard removes all code logic from the original tasks, retaining only file dependency relationships to ensure automated scripts remain functional.
Agents receive \emph{no compiler construction knowledge} and \emph{no implementation framework}---they must implement the full compiler from scratch.
The task descriptions are reduced to experiment requirements and scoring script usage instructions, with all code logic explanations and framework references removed.

This workload measures an agent's ability to \emph{construct} a complete system from minimal specification, representing a fundamentally harder evaluation setting.
The difficulty gap between \yatcc{} and \yatcc{}-Hard quantifies the value of scaffolding in agentic coding tasks.

\subsection{Resurrection Protocol}

The \resurrection{} mechanism addresses the fundamental evaluation challenge of serial workloads: when an agent fails at Task~$k$, downstream tasks become unreachable.
\resurrection{} injects the \emph{golden intermediate artifact} (the correct output from the reference implementation) after failure, allowing the agent to continue to Task~$k+1$ with corrected state.

\textbf{Definition.}
When Task~$k$ scores below the passing threshold (60\%), \resurrection{} is triggered:
(1)~the agent's failed artifact is replaced with the golden artifact;
(2)~the build system is reconfigured via \texttt{config.cmake} (\texttt{TASK\_k\_REVIVE=ON});
(3)~the agent continues to Task~$k+1$ with the corrected state.

\textbf{Design rationale.}
Resurrection is an \emph{evaluation protocol}, not a scoring adjustment.
It decomposes the compound failure ``cannot reach Task~$k+1$'' into two independent measurements:
(1)~\emph{predecessor failure}: the agent could not produce a passing artifact at Task~$k$;
(2)~\emph{downstream capability}: given correct prerequisites, can the agent succeed at Task~$k+1$?

This decomposition enables node-level pass rate reporting even when the full pipeline is broken, significantly improving diagnostic resolution.
Agents are \emph{not informed} about resurrection---the mechanism is triggered externally by the evaluation orchestrator, eliminating moral hazard.

\subsection{Intermediate Artifacts and Verification}

Each task produces a well-defined intermediate artifact:
Task~1 produces a \emph{token stream}; Task~2 produces an \emph{AST/ASG}; Task~3 produces \emph{LLVM IR}; Task~4 produces \emph{optimized IR}; Task~5 produces \emph{RV64 assembly}.
Scoring is performed by dedicated test suites that execute the agent's code against comprehensive test cases.
Each task's score ranges from 0 to 100, with a passing threshold of 60\%.